\documentclass{article}
\usepackage[utf8]{inputenc}
\usepackage{amsmath}
\usepackage{amsfonts}
\usepackage{amssymb}
\usepackage{geometry}
\usepackage{tcolorbox}
\usepackage{hyperref}
\geometry{a4paper, margin=1in}

\title{Module 08: Production Pipelines \& Collaborative Git Workflows - Part 3}
\author{Cybernauts 2026 Datathon Campaign}
\date{May 2026}

\begin{document}

\maketitle

\begin{abstract}
A brilliant machine learning pipeline is worthless if your team cannot integrate it under hackathon pressure. When multiple data scientists collaborate on a shared codebase simultaneously, uncoordinated pushes will break execution paths and destroy valuable work. This note establishes the high-velocity Git development strategies and conflict resolution mechanics required to maintain an uncorrupted repository during a fast-paced datathon sprint.
\end{abstract}

\tableofcontents
\newpage

\section{The Chaos of Direct Collaboration}
During the high-stress environment of a hackathon, speed is essential. However, sacrificing code management for raw speed results in an engineering bottleneck.

\subsection{The Main Branch Vulnerability}
Allowing multiple team members to push unverified experimental code directly to the \texttt{main} branch is an anti-pattern. If Team Member A pushes an un-vetted feature that contains a syntax error while Team Member B is in the middle of a model tuning run, Team Member B's environment will instantly crash. The repository becomes broken, and valuable development cycles are wasted debugging environment state failures rather than refining models.

\section{The Branch-per-Feature Strategy}
To maintain continuous development alignment, the codebase must treat the master line as a sacred space.

\subsection{The Pristine Production Standard}
The \texttt{main} branch must remain clean and production-ready at all times. It should house a completely functional, end-to-end baseline pipeline that can be executed cleanly from start to finish via \texttt{python main.py} without throwing exceptions.

\subsection{Feature-Driven Isolation}
Every single experiment—whether it is a new data transformation, an alternative model exploration, or an optimization script—must be developed in isolation on a dedicated branch spun off from \texttt{main}. This is implemented using explicit branch naming conventions:
\begin{verbatim}
git checkout -b feature/transaction-aggregations
git checkout -b model/optuna-lgbm-tuning
\end{verbatim}

By isolating development streams, teammates can build out features or run heavy hyperparameter searches in parallel without affecting the core baseline.

\section{Atomic Commits and Granular History}
When writing code within isolated branches, avoid the trap of bundling hundreds of lines of varied edits into a single, generic save point like \texttt{git commit -m "fixed bugs"}.

\subsection{The Logic of Single-Intent Changes}
Adhere to the principle of \textbf{Atomic Commits}. Each commit must represent a single, isolated, logical unit of work, accompanied by an explicit description:
\begin{verbatim}
git commit -m "feat: add user spend transform feature"
\end{verbatim}

\textit{The Competitive Safety Net:} If Team Member A integrates their feature branch into the production pipeline, but the collective cross-validation score unexpectedly drops, atomic logs allow the team to pinpoint the exact commit that degraded performance. The team can cleanly roll back the single broken feature script without losing any other improvements made by the rest of the team.

\section{Mastering Merge Conflict Resolution}
When multiple developers modify the same file simultaneously, Git cannot automatically reconcile the changes. This triggers a merge conflict.

\subsection{Deconstructing Git Conflict Markers}
When a conflict occurs in a core script (e.g., \texttt{src/features.py}), Git halts the merge and injects structural boundary markers into the affected file:
\begin{verbatim}
<<<<<<< HEAD
    # Current Branch: Team Member A's implementation
    df['user_income_log'] = np.log1p(df['income'])
=======
    # Incoming Branch: Team Member B's implementation
    df['income_per_capita'] = df['income'] / df['family_size']
>>>>>>> feature/income-transforms
\end{verbatim}

To resolve the conflict, you must understand these structural boundaries:
\begin{itemize}
    \item \texttt{<<<<<<< HEAD}: Marks the start of the conflicting code on your active branch.
    \item \texttt{=======}: The dividing baseline that separates the two conflicting code segments.
    \item \texttt{>>>>>>> branch\_name}: Marks the end of the conflicting code from the incoming branch.
\end{itemize}

\subsection{The Manual Resolution Protocol}
To fix the conflict, developers must open the file, evaluate both implementations, and manually edit the code block. You must determine whether to keep one option, reject both, or blend them into a unified structure. Finally, remove all injected conflict markers (\texttt{<<<}, \texttt{===}, \texttt{>>>}) and finalize the merge with a clean validation commit:
\begin{verbatim}
git add src/features.py
git commit -m "fix: resolve merge conflict in feature script"
\end{verbatim}

\newpage
\section{Practice Problems}

\textbf{P1:} A team member finishes writing an automated feature selection script on branch \texttt{feature/selection}. They want to merge their work into the production \texttt{main} branch. Write down the sequence of terminal Git commands required to safely pull the latest updates from the remote repository, move to the target branch, and execute the merge.
\\
\textit{Answer:}
\begin{verbatim}
# Fetch and pull latest pristine remote updates
git checkout main
git pull origin main

# Switch to the experimental feature branch and rebase/merge to incorporate main's updates
git checkout feature/selection
git merge main

# Switch back to main and safely merge the updated feature branch
git checkout main
git merge feature/selection
git push origin main
\end{verbatim}

\newline
\textbf{P2:} Two data scientists are editing the global parameter file \texttt{src/config.py} in parallel. When merging, Git halts with a conflict marker indicating overlapping edits in the hyperparameter dictionary. Team Member A changed the learning rate to $0.01$, while Team Member B changed it to $0.03$. Explain the step-by-step process required to resolve this conflict without breaking the pipeline.
\\
\textit{Answer: First, open the conflicted \texttt{src/config.py} file in a text editor and locate the conflict markers. The team must communicate to decide which learning rate parameter performed better within their cross-validation loop. Once chosen (e.g., $0.03$), delete Team Member A's code block along with the \texttt{<<<<<<< HEAD}, \texttt{=======}, and \texttt{>>>>>>>} markers, leaving only the chosen configuration line. Save the file, run a quick baseline evaluation to verify functionality, and execute \texttt{git add} and \texttt{git commit} to finalize the merge.}

\newline
\textbf{P3:} Why does practicing atomic commits significantly lower validation risks when integrating multiple complex feature transformations into a shared team codebase?
\\
\textit{Answer: If multiple features are bundled into a single commit, and the final cross-validation score drops, it is difficult to determine which specific transformation introduced the issue. By using atomic commits, each feature is isolated in the git history. If an issue arises, the team can use granular debugging tools (like \texttt{git bisect} or \texttt{git revert}) to identify and remove the single problematic transformation while keeping all other high-performing features intact.}

\end{document}
